% Transcription — fonds Alexandre Grothendieck, shelfmark 29, batch 10 (pages 181-200).
% Established from the facsimile archives/batches/29/batch-10.pdf (a mirror of
% https://grothendieck.umontpellier.fr/29.pdf), per the skill
% transcribe-grothendieck.
%
% Pass: Opus 5 (claude-opus-5), 2026-08-29 — first pass, unchecked against the
% pages by a human.
%
% Eighteen of the twenty pages are transcribed. Two are not:
%
%   - page 188, a brown cardboard wrapper inscribed top right, in his hand,
%     « Domaines de ramification / Catégories multigaloisiennes »;
%   - page 189, an otherwise blank leaf.
%
% The wrapper names the run that follows it, exactly as pages 111, 123, 143,
% 163, 172, 174 and 177 of the earlier batches; it gets no \page{}, and its
% words are carried into the section heading instead, where they are marked as
% his.
%
% The batch falls into four runs.
%
%   Pages 181-185 — the continuation of batch 9's last computation: the
%     semi-cosimplicial functor F⁰ ⇉ F¹ ⇛ F², the functors π⁰(F*,ξ;E) and
%     π¹(F*,ξ;G) attached to it, and strict pro-representability. Pages 181,
%     182, 184 and 185 are pale pencil on a leaf whose verso shows through at
%     nearly the same tone — 181 carries a full ghost of 182 — so the two sides
%     have to be separated before anything can be read. The formulas carry; a
%     good deal of the connective prose does not.
%   Page 183 — a landscape sheet of sketches: an edge path from x to a, the
%     coboundary of a 1-cochain, π₁^#(S,a), k[[s,t]]. It has no running
%     argument and the transcription does not impose one.
%   Page 186 — a proposition, in blue ink and written out fair, on when a
%     morphism is of finite type.
%   Page 187 — « Opérateur d'homotopie pour d'' », cancelled by two long
%     diagonals. It reads end to end under them.
%   Pages 190-200 — the run the wrapper of page 188 names, and the substance of
%     the batch: domaines de ramification, their axioms R1-R4, the topology
%     they carry, Prop. 1 to Prop. 6, the 2-category Domram(S), and the
%     Kummer examples. He paginates it 1-10 in the top right corner; page 193
%     carries no number of his and is an inserted proposition. His numbering is
%     recorded at each page where it is visible. The run breaks off mid-sentence
%     at page 200 and continues into batch 11.
%
% Two notations recur and are fixed once for the batch. He writes ½ for
% "semi-": « ½cosimplicial » is semi-cosimplicial, and his s can be read as a 2,
% so the word looks like "½co2-indicat" until it is read against the diagram
% F⁰ ⇉ F¹ ⇛ F² that it names. And he underlines a family to mean the whole
% family: a̲ = (a_i), n̲ = (n_i), so Z^a̲_n̲ is the Kummer covering they define.
%
% The dating is the inventory's own, for the whole shelfmark.
\documentclass[11pt,a4paper]{article}
% Le préambule commun à toutes les transcriptions du fonds Grothendieck.
%
% Il tient en une idée : **l'incertitude se compose, elle ne se résout pas.**
% Une transcription qui lisse un mot douteux a détruit la seule chose pour
% laquelle on la faisait — un lecteur doit pouvoir savoir, à chaque endroit,
% ce qui était sur le feuillet et ce qui a été ajouté. Les six macros qui
% suivent sont donc l'appareil critique, et non de la décoration.
%
% Les mêmes commandes sont reconnues par `scripts/render.mjs`, qui produit la
% vue de lecture du site. Une macro ajoutée ici doit l'être aussi là-bas,
% faute de quoi le rendu échouera bruyamment — ce qui est voulu : un rendu
% silencieusement faux est pire qu'un rendu absent.

\ProvidesPackage{grothendieck}[2026/08/08 Transcriptions du fonds Grothendieck]

\usepackage{amsmath,amssymb,amsthm}
\usepackage{mathrsfs}
\usepackage{stmaryrd}
% Les diagrammes commutatifs sont partout dans ce fonds : c'est la langue
% même de la géométrie algébrique de Grothendieck, et une transcription qui
% les rendrait en prose ne transcrirait rien.
\usepackage{tikz-cd}
% Français par défaut — c'est la langue des feuillets — mais l'anglais est
% chargé aussi : la traduction et le résumé sont des documents anglais, et la
% typographie française (espaces avant les deux-points) y serait une faute.
% Ces documents basculent par \selectlanguage{english} en tête de corps.
\usepackage[english,french]{babel}
\usepackage{fontspec}
\usepackage{geometry}
\geometry{a4paper,margin=25mm,marginparwidth=28mm}
\usepackage{marginnote}
\usepackage{xcolor}
% ulem, not soul. `soul`'s \st and \ul are built on a character-by-character
% scan that XeLaTeX's font machinery breaks: the document fails with an
% undefined control sequence rather than degrading. ulem does the same two jobs
% and is Unicode-engine safe. `normalem` stops it hijacking \emph, which the
% transcriptions use for Grothendieck's own emphasis.
\usepackage[normalem]{ulem}
\usepackage{hyperref}
\hypersetup{colorlinks=true,linkcolor=black,urlcolor=black,pdfborder={0 0 0}}

% --- Métadonnées du lot -----------------------------------------------------
% Reprises telles quelles par le script de rendu, qui les affiche en tête de
% la vue de lecture. Elles fixent ce que le fichier prétend couvrir : sans
% elles, une transcription flotte sans point d'attache dans un fonds de seize
% mille feuillets.

\newcommand{\folder}[1]{\gdef\grothFolder{#1}}
\newcommand{\batch}[1]{\gdef\grothBatch{#1}}
\newcommand{\leaves}[2]{\gdef\grothFirst{#1}\gdef\grothLast{#2}}
\newcommand{\foldertitle}[1]{\gdef\grothFoldertitle{#1}}
\gdef\grothFolder{?}\gdef\grothBatch{?}\gdef\grothFirst{?}\gdef\grothLast{?}\gdef\grothFoldertitle{}

% --- Le résumé --------------------------------------------------------------
%
% Il ouvre la lecture modernisée, sous le titre « Résumé », et il est composé
% en petit. Ce n'est pas un ornement : c'est le seul passage du document qui
% ne soit pas des mathématiques, et un lecteur doit pouvoir le reconnaître
% avant de l'avoir lu — puis le sauter s'il connaît déjà le sujet, ou s'y
% arrêter s'il ne le connaît pas. Le filet qui le suit marque la reprise du
% corps.

\newenvironment{resume}
  {\small\setlength{\parskip}{0.45em}}
  {\par\medskip\hrule\medskip}

% --- Mots-clés ---------------------------------------------------------------
%
% En anglais, en clôture du résumé de la lecture modernisée : le vocabulaire
% moderne sous lequel chercher ce que les feuillets construisent. Le manifeste
% du site (`npm run manifest`) les extrait pour étiqueter la cote entière —
% cette ligne est la source unique des tags, il n'y a pas de fichier à côté.
\newcommand{\keywords}[1]{\par\smallskip\noindent
  {\footnotesize\textsc{Keywords} --- \textit{#1}}\par}

% --- L'appareil critique ----------------------------------------------------

% Le numéro de feuillet, en marge. C'est l'ancre qui permet de régler un
% désaccord sur une formule en regardant *un* feuillet plutôt qu'une chemise
% de six cents. Le site s'en sert aussi pour tourner les pages du fac-similé
% quand on fait défiler la transcription.
\newcommand{\leaf}[1]{%
  \marginnote{\footnotesize\color{gray!70}\textsf{#1}}%
  \label{leaf:#1}\ignorespaces}

% Illisible : on ne devine pas. Un mot inventé qui se lit comme les autres est
% le pire résultat possible.
\newcommand{\ill}{\textcolor{red!60!black}{[\,\ldots\,]}}

% Lecture douteuse : le mot est donné, et signalé comme incertain.
\newcommand{\uncertain}[1]{\uline{#1}}

% Ajout de l'éditeur — une lettre manquante, un mot sous-entendu.
\newcommand{\add}[1]{\textcolor{blue!50!black}{[#1]}}

% Biffé par Grothendieck lui-même. Ce qu'il a rayé fait partie du document :
% une rature dit souvent où la pensée a changé de direction.
\newcommand{\struck}[1]{\sout{#1}}

% Note du transcripteur, sur l'état matériel du feuillet.
\newcommand{\note}[1]{\par\smallskip{\small\color{gray!60!black}\textit{[#1]}}\par\smallskip}

% Note marginale de Grothendieck, distincte du corps du texte.
\newcommand{\marginal}[1]{\par\smallskip\noindent\hspace{1em}%
  {\small\color{gray!60!black}#1}\par\smallskip}

% --- Titre ------------------------------------------------------------------

\AtBeginDocument{%
  \begin{center}
    {\large\bfseries Fonds Alexandre Grothendieck --- cote n\textsuperscript{o}~\grothFolder}\\[2pt]
    {\itshape\grothFoldertitle}\\[4pt]
    {\small feuillets \grothFirst--\grothLast\ (lot \grothBatch)}\\[2pt]
    {\footnotesize\color{gray!60!black}Transcription établie d'après le fac-similé numérisé
     par l'Université de Montpellier.\\
     Les lectures incertaines sont soulignées, les illisibles notées [\,\ldots\,],
     les ajouts entre crochets.}
  \end{center}
  \vspace{1em}}

\endinput


\folder{29}
\batch{10}
\pages{181}{200}
\dating{1959-1969}
\watermark{Édition de démonstration}
\foldertitle{Groupe fondamental [Autour de SGA 4 et SGA 7] : lettres (1959, 1967, 1969), tapuscrits et copies de tapuscrit annotés (s.d.), notes manuscrites (s.d.)}

\begin{document}

\section*{$\pi_0$ : foncteurs ½cosimpliciaux, pro-représentabilité stricte (pages 181, 182, 184 et 185)}

\note{le titre « $\pi_0$ » est de lui et ouvre la page 182 ; les pages 181, 182,
184 et 185 sont d'un crayon très pâle sur un feuillet dont le verso transparaît
au même ton, et la page 181 porte en filigrane la page 182 tout entière ---
seule l'encre de chaque recto est transcrite ici}

\page{181}

Pratiquement, \uncertain{on} vérifie : \uncertain{contractions}
$i_0(k_0)$, $i_1(k_1)$, $i_2(k_2)$ \ill{} et \ill{} de bonnes
\uncertain{cochaînes} \ill{}, sur les \struck{\ill{}} \ill{}.

\note{les trois premières lignes de la page sont barrées en leur milieu par une
insertion elle-même biffée ; ce qui subsiste ne se laisse pas lire}

\begin{tikzcd}[column sep=small, row sep=small]
\overline{K}_0 & \overline{K}_1 \arrow[l, bend left=12] \arrow[l, bend right=12]
  & \overline{K}_2 \arrow[l, bend left=20] \arrow[l] \arrow[l, bend right=20]
\end{tikzcd}

On utilise le fait que
\[
  \pi_0(k_{\ast}) = e, \qquad \text{« } \pi_1(k_{\ast}) = e \text{ »}
  \qquad \text{i.e.} \quad \pi^1(k_{\ast}, G) = e
\]

\note{l'indice de $k$ est un pâté, comme au bas de la page 180 ; le contexte ne
le fixe pas et il est rendu ici par $\ast$}

donc t\add{ou}te donnée de descente est une \uncertain{prolongation} et on
utilise le fait
\[
  \pi_0(K_{\ast}) = e
\]
pour le fait \add{que} \struck{deux} \struck{prolongations} d'un \ill{}
\uncertain{automorphisme} d'un \ill{} structure de descente \ill{}
\struck{pro-}\ill{} est trivial.

\note{suit un bloc encadré, relié par des arcs à ce qui précède, et dont la
lecture ne se laisse pas établir ; il se termine sur un mot biffé,
\struck{\uncertain{Régularisation}}}

\marginal{\ill{}}

\page{182}

$\pi_0$

\begin{tikzcd}[column sep=small, row sep=small]
F^0 \arrow[r, bend left=12] \arrow[r, bend right=12]
  & F^1 \arrow[r, bend left=20] \arrow[r] \arrow[r, bend right=20] & F^2
\end{tikzcd}

foncteurs $(\mathrm{Ens}) \to (\mathrm{Ens})$, objet simplicial de
$\mathrm{Hom}(\mathrm{Ens}, \mathrm{Ens})$.

\uncertain{Posons} en principe \ill{} définir les $\pi_0$, $\pi_1$.

$\pi_0(F)$, $F \in \underline{\mathrm{Hom}}(\mathrm{Ens}, \mathrm{Ens})$, $F$
\[
  F(I) \underset{?}{\simeq} C(K, I)
\]
\marginal{$K$ un gpe top. tot. discontinu [\ill{}]}

\marginal{cond. néc.}
\begin{enumerate}
\item[a)] $F$ commute aux \struck{produits} $\varprojlim$ finies, i.e. $F$ exact
à gauche [\ill{} $F(\emptyset)$ \ill{}]
\item[b)] $F(\bigcap I_{\alpha}) = \bigcap_{\alpha} F(I_{\alpha})$ pour
\struck{toute} famille \uncertain{filtrante} de sous-ensembles de $I$
\item[c)] \struck{Alors} $\varprojlim \pi_{0j} \longrightarrow \pi_{0i}$ est
\emph{surjectif} pour tt $i$
\end{enumerate}

Ces conditions équivalent à : $F$ \emph{strictement pro-représentable} par un
système projectif strict $(\pi_{0i})_i$ \ill{} d'ensembles.

Ces conditions sont \emph{néc.} et \emph{suffisantes}, et on pourra prendre
$K = \varprojlim \pi_{0i}$, avec la topologie limite inverse des top. discrètes.

\emph{N.B.} \ill{} $F$ est \ill{} de $2$ \uncertain{façons} \add{pointées}
$C(L, I)$, $L$ un gp. top., \ill{} $\pi_0(L) \to K$. Si $L$ n'est
\uncertain{pas} compact, ce \ill{} n'est \ill{} un \uncertain{isomorphisme}, ni
\uncertain{surjectif}, ni \uncertain{injectif},

\marginal{Bourb. Topologie \uncertain{II}, p. 232, prop. 6}

\note{le chiffre romain du fascicule est de ceux que sa main ne distingue pas
d'un 2 ; il est laissé tel quel}

\section*{Feuillet d'esquisses (page 183)}

\note{page en travers, couverte de tracés sans ordre de lecture ; on en donne
ci-dessous les fragments qui se lisent, dans leur disposition sur la feuille,
sans leur supposer d'enchaînement}

\page{183}

\begin{tikzcd}[column sep=small, row sep=small]
\overline{K}_0 & \overline{K}_1 \arrow[l, bend left=12] \arrow[l, bend right=12]
  & \overline{K}_2 \arrow[l, bend left=20] \arrow[l] \arrow[l, bend right=20]
\end{tikzcd}

\note{deux traits de renvoi descendent de ce diagramme vers les mots
\emph{fini} et \emph{gpe top. discret}}

$\Lambda$ \qquad $\mathbb{Z}$ \qquad $F$ \qquad $f(x,y)$ \qquad $f(u)$
\qquad $K_1$ \qquad $E$ \qquad $S \leftarrow S'$

$\overline{K}_1 \to G$ qui \ill{} un $\overline{F}$ \ill{} de façon que
$\overline{K}_1$ y \ill{} \emph{continue}

$k[[s,t]]$, \quad $k$ alg. clos

\[
  g(p(t_1))\, f(t_1)\, g(q(t_1))^{-1}
\]

\note{les arguments des trois facteurs sont écrits menu et ne se distinguent
pas sûrement les uns des autres ; la formule est celle de l'équivalence des
données de descente de la page 180}

$\pi_1^{\#}(S, a)$, encadré, et surmonté de \emph{$S$ connexe}

\note{le signe entre $\pi_1$ et $(S,a)$ est un croisillon, tracé deux fois sur
la page de la même façon ; c'est donc une notation et non une rature, mais rien
ici n'en fixe le sens}

\note{au bas de la feuille, un chemin d'arêtes de $x$ à $a$, dont les flèches
ne vont pas toutes dans le même sens :
$x \xrightarrow{u_n} \cdot \xrightarrow{u_{n-1}} \cdot \; \cdots \;
\xrightarrow{u_1} a$ ; à côté, $d_0$, $d_1$, $d_2$ portés sur une esquisse
simpliciale}

\section*{Le foncteur noyau $\pi^0(F^{*})$, et les $\pi^0(F^{*},\xi;E)$, $\pi^1(F^{*},\xi;G)$ (pages 184 et 185)}

\page{184}

mais il est naturel de mettre sur $\pi_0(L)$ la topologie image inverse de
\add{celle de} $K$ \add{topologie \struck{\ill{}} des $\pi_0(L)$},
pratiquement, \ill{} de la\,]. $F$ est \ill{} \uncertain{connexe} si
$F \simeq \mathrm{id}_{\mathrm{Ens}}$, i.e. $F$ est $T$-représentable par
l'objet ponctuel [\ill{} par une \ill{} \ill{}].

Considérons \struck{un} deux foncteurs \ill{}
\[
  F^{*} : \quad F^{*0} \rightrightarrows F^{*1}
  \qquad (F^0, F^1 \in \underline{\mathrm{Hom}}(\mathrm{Ens}, \mathrm{Ens}))
\]
et le \struck{foncteur} \emph{noyau}
\[
  \pi^0(F^{*}) = F = \mathrm{Ker}\,(F^0 \rightrightarrows F^1)
\]
si $F^0$, $F^1$ sont \struck{strictement} \add{$T$-}\,pro-représentables, alors
les deux flèches précédentes \ill{} des flèches

\begin{tikzcd}[column sep=small, row sep=small]
K_0 & K_1 \arrow[l, bend left=12] \arrow[l, bend right=12]
\end{tikzcd}

et $F$ sera également $T$-représentable [\emph{N.B.} Peut-être \ill{} des
foncteurs $T$-représentables ou $T$-\uncertain{groupes}, \ill{} prennent \ill{}
limites inductives \ill{} \ill{} pro-\ill{} et le $\pi_0$ associé\,]. Une \ill{}
où $F^{*}$ \ill{} donnée est \ill{} et \ill{} son noyau est
\uncertain{connexe}.

Considérons un « foncteur ½cosimplicial », \add{de}
$\underline{\mathrm{Hom}}(\mathrm{Ens}, \mathrm{Ens})$

\begin{tikzcd}[column sep=small, row sep=small]
F^0 \arrow[r, bend left=12] \arrow[r, bend right=12]
  & F^1 \arrow[r, bend left=20] \arrow[r] \arrow[r, bend right=20] & F^2
\end{tikzcd}

on lui associe \ill{} foncteurs $\pi^0(F^{*}; E)$, $\pi^1(F^{*}; G)$, si $E$ un
ensemble, $G$ un groupe variable. Si $F^{*} \to \mathrm{id}_{\mathrm{Ens}}$, on
définit \struck{\ill{}} i.e. muni de $\pi^0(F^{*}, \xi; E)$ et
$\pi^1(F^{*}, \xi; G)$,

\page{185}

désignant la \emph{ponctuation}. Plus précisément, si on a \ill{} un \ill{}
foncteur ½cosimplicial
\[
  f^0 \rightrightarrows f^1 \Rrightarrow f^2
\]
\struck{\ill{}} et un \uncertain{lien} ½cosimplicial $\xi$

\begin{tikzcd}[column sep=small, row sep=small, nodes={font=\scriptsize}]
F^0 \arrow[r, bend left=12] \arrow[r, bend right=12] \arrow[d, "\xi^0"]
  & F^1 \arrow[r, bend left=20] \arrow[r] \arrow[r, bend right=20]
    \arrow[d, "\xi^1"] & F^2 \arrow[d, "\xi^2"] \\
f^0 \arrow[r, bend left=12] \arrow[r, bend right=12]
  & f^1 \arrow[r, bend left=20] \arrow[r] \arrow[r, bend right=20] & f^2
\end{tikzcd}

on définit $\pi_0(F^{*}, \xi; E)$ et $\pi^1(F^{*}, \xi; E)$.

Si $F^0$, $F^1$, $F^2$ sont $T$-représentables, et dans le cas \ill{} où
\ill{} le \ill{} $\xi : F^{*} \to f^{*}$, si $f^{*}$ est $T$-représentable, on
ramène ceci à la donnée d'une \struck{\ill{} cosimplicial}

\begin{tikzcd}[column sep=small, row sep=small]
K_0 & K_1 \arrow[l, bend left=12] \arrow[l, bend right=12]
  & K_2 \arrow[l, bend left=20] \arrow[l] \arrow[l, bend right=20]
\end{tikzcd}

\ill{} d'un \emph{\uncertain{ensemble} ½simplicial ponctué} (\ill{})

\begin{tikzcd}[column sep=small, row sep=small]
k_0 & k_1 \arrow[l, bend left=12] \arrow[l, bend right=12]
  & k_2 \arrow[l, bend left=20] \arrow[l] \arrow[l, bend right=20]
\end{tikzcd}

par un procédé \emph{standard}\,\ldots

Il y a lieu de généraliser ces constructions --- on a vu que des $K_i$,
\struck{\ill{}} des $\overline{K}_i$, on a des catégories galoisiennes
\struck{itérées} \ill{} à chaque pt, (\struck{\ill{}, les $F_{\xi}$}) pour
\add{dé}cider \ill{} de descente : le \emph{Van Kampen}\,\ldots

\section*{Une proposition sur les morphismes de type fini (page 186)}

\page{186}

\textbf{Proposition.} Soit $f : X \to Y$ un morphisme, avec
\begin{enumerate}
\item[1)] \add{$Y$ loc. noeth. et \ill{} de type fini}
\struck{$f$ loc. de présentation finie et \ill{}}
\item[2)] $X$ \struck{irréductible} \ill{} \add{n'a qu'un} \ill{} fini de comp.
irréd.
\item[3)] $f$ loc. \add{de type fini} \struck{séparé-fini} \struck{et} $f$
\emph{séparé}
\end{enumerate}

Sous ces conditions, $f$ est \emph{de type fini}.

On peut supposer \struck{$X$} $Y$ affine, donc noethérien, \ill{} (car il
suffit de voir pour $f^{-1}$ \add{d'un} pt. générique). $X$ intègre
\struck{et $f$ dominant}. Soient $U$ un \ill{} de $X$, il est \ill{} au-dessus
de $Y$, dont \ill{} mais \ill{} il est \uncertain{½ isomorphe} \ill{} fini
$X'$, qui \ill{} un \ill{} intègre. Soit $X''$ \ill{}

\begin{tikzcd}[column sep=small, row sep=small, nodes={font=\scriptsize}]
& X'' \arrow[dl, "\text{fini}"'] \arrow[d] \\
X \arrow[d] & X' \\
Y &
\end{tikzcd}

\note{le diagramme est porté dans la marge de gauche ; la flèche $X'' \to X$
est annotée \emph{fini}, celle qui descend vers $Y$ porte une mention
illisible}

\ill{} fini intègre de $X \times_Y X'$ \struck{\uncertain{adhérence}} \ill{} le
« \ill{} » de $U \times_Y U'$, dans $X''$ ; ainsi \ill{} \add{au-dessus de} $X$
et $X'$.

\note{le mot entre guillemets nomme la construction et ne se laisse pas lire ;
il n'est pas remplacé par une conjecture}

Mais $X'' \to X$ est \emph{fini}, et il est dominant, \struck{est}
\emph{surjectif}. Il suffit donc de voir que $X''$ est de t.f. sur $Y$, on
\ill{} \add{remplacer $X$ par} $X'$. Or $X'' \to X'$ est un morphisme
\emph{birationnel}, \ill{} $X''$ et $X'$ intègres. Donc \struck{il y a} un
$X'$-morphisme \struck{surjectif} $X' \to X''$, i.e. $X'$ est le
\uncertain{normalisé}. Or $X' \to$ \ill{} fini, \ill{} \ill{}\,\ldots

\section*{« Opérateur d'homotopie pour $d''$ » (page 187)}

\note{la page est annulée par deux longues diagonales ; elle se lit néanmoins
d'un bout à l'autre, et la biffure n'est signalée qu'ici}

\page{187}

\emph{Opérateur d'homotopie pour $d''$}

\[
  \overline{W}_n(G) = G_{n-1} \times G_{n-2} \times \cdots \times G_0
\]
\struck{$\overline{W}_n(G)$} \quad $\tau : \overline{W}_n(G) \longrightarrow G_{n-1}$

\[
  \varphi_n(x_1, \ldots, x_n) = \frac{1}{n}\left(1 + \frac{p_1}{n-1}
  + \frac{p_2}{\frac{(n-1)(n-2)}{2}} + \cdots + \frac{p_k}{\binom{n-1}{k}}
  + \cdots + p_{n-1}\right)
\]
où $p_k = \struck{\text{somme des}} \sum x_1 x_2 \ldots x_k$.

\note{les indices des variables sommées sont écrits sans distinction ; la somme
est celle des fonctions symétriques élémentaires}

\[
  U_j f = \frac{1}{2\pi i} \int_{B_j}
    f(z_1, \ldots, z_{j-1}, t, z_{j+1}, \ldots, z_n)\, \frac{dt}{t - z_j}
\]
\[
  T_j f = \frac{1}{2\pi i} \iint_{K_j}
    f(z_1, \ldots, z_{j-1}, t, z_{j+1}, \ldots, z_n)\,
    \frac{dt\, \overline{dt}}{t - z_j}
\]

\note{la variable complexe est tracée en 3 bouclé, comme le groupe sans nom de
la page 176 ; l'indice porte le même trait, et il est rendu ici par $j$ puisque
le contexte en fait un entier entre $1$ et $n$}

\struck{Et} Pour \struck{toute} suite strictement croissante
$j_1 < j_2 < \cdots < j_q$, on \struck{note} introduit l'opérateur
$\varphi_n(U ; j_1, \ldots, j_q)$ \struck{\ill{}} ce que devient
$\varphi_n(U_1, \ldots, U_n)$ quand on remplace $U_{j_1}, \ldots$ et $U_{j_q}$
par \emph{zéro}.

L'opérateur d'homotopie, pour
$\omega = a\, d\overline{z}_{j_1} \wedge \cdots \wedge d\overline{z}_{j_q}$
$(q \geq 1)$ est
\[
  P\omega = \varphi_n(U ; j_1, \ldots, j_q) \ast \sum_{\gamma=1}^{q}
  (-1)^{\gamma-1} (T_{j_\gamma} a)\;
  d\overline{z}_{j_1} \wedge \cdots \wedge
  \widehat{d\overline{z}_{j_\gamma}} \wedge \cdots \wedge
  d\overline{z}_{j_q}
\]
\marginal{($n$ est le nombre des variables)}

\note{le signe qui sépare $\varphi_n$ de la somme est un astérisque appuyé ;
la page ne dit pas de quelle opération il s'agit}

\[
  P d'' f \;\struck{d'' P f} \;=\; f - U_1 U_2 \ldots U_n f
\]

\section*{Domaines de ramification, catégories multigaloisiennes : les axiomes R1 à R4 (page 190)}

\note{le feuillet qui ouvre la suite --- page 188, une chemise de carton par
ailleurs nue --- porte de sa main, souligné deux fois, « Domaines de
ramification / Catégories multigaloisiennes » ; ces mots donnent son titre à la
suite et le feuillet n'est pas transcrit. La page 189 est nue. Il pagine
lui-même la suite de 1 à 10, en haut à droite}

\page{190}\note{sa page 1}

$\mathrm{Et}(S)$ le site étale de $S$. $\mathcal{U}$ un univers.

$S$ schéma. Un « domaine de ramification » sur $S$ est un couple $(R, r)$ formé
d'un champ \struck{en} \add{$\mathcal{U}$-}\emph{catégories multigaloisiennes}
$R$ sur $\mathrm{Et}(S)$, et d'un morphisme \struck{de champs}
\add{\emph{continu}} (« réalisation géométrique »)
\[
  r : R \longrightarrow \mathrm{Fl}_S
\]
où $\mathrm{Fl}_S$ est la \struck{catégorie fibrée} \ill{} sur
$\mathrm{Et}(S)$ dont la fibre en $S'$ est $(\mathrm{Sch})_{/S'}$ ;
\struck{\ill{}} un objet $X$ de $R_{S'}$ (\struck{objet} de $\mathrm{Et}(S)$)
est dit \emph{revêtement de type $R$ sur $S'$}, \ill{} $r(X)$ est sa
\emph{réalisation géométrique}. On \struck{fait l'hypothèse} \ill{} les données
sont soumises aux axiomes suivants :

\begin{itemize}
\item[\textbf{R1}] \struck{Le foncteur} $r$ commute aux \emph{sommes finies} et
\struck{au passage au quotient} par une relation d'équivalence [il serait
$=$ axiome \struck{identique}, $\nexists$ \emph{plus bas} \ill{}, mais en
général pas exact à gauche !]. (\ill{} il suffit \ill{} passage au quotient
\ill{} \ill{} catégories galoisiennes \ill{})

\item[\textbf{R2}] \struck{Donné} Soient $S' \in \mathrm{Ob}\,\mathrm{Et}(S)$,
et $X \in \mathrm{Ob}\,R_{S'}$, l'application $X' \mapsto r(X')$
(\add{$=$ ss.-objets}) est une \emph{bijection} de l'ens. des
\struck{sous-objets} directs $X'$ de $X$ avec l'ens. des sous-ensembles
directs (i.e. parties ouvertes et fermées) de $r(X)$.

\item[\textbf{R3}] Soient $S'$, $X$ comme dessus, et soit \struck{\ill{}}
$(X_i)$ une famille (\uncertain{pas néc.} finie) de \struck{sous-objets de $X$},
telle que la famille des $r(X_i)$ recouvre $r(X)$. Alors la famille
$X_i \to X$ est une famille de descente \emph{effective}
($=$ \struck{un épimorphisme} \ill{} effectif \ill{}).
$\gamma = \coprod_{i \in I} X_i$, avec $X \times_Y X \simeq{}$ \ill{}

\item[\textbf{R4}] Soit $X \to Y$ dans $R_S$, \ill{}
$S' \in \mathrm{Ob}\,\mathrm{Et}(S)$, \ill{} dominant. On veut que les $r(X_i)$
recouvrent $r(Y)$. \ill{} $(X_i \to X)$
\end{itemize}

\textbf{Prop. 1} \struck{Écri} Disons qu'une famille de flèches de $R_{S'}$ est
\emph{convergente} si la famille $r(X_i) \to r(X)$ est \emph{surjective}, i.e.
(grâce à R1 c'est « Kif Kif ») et les images $X_i \subset X$ sont telles que les
$r(X_i)$ recouvrent $r(X)$. Alors \struck{\ill{}} ces \emph{familles} sont les
familles couvrantes d'une \emph{topologie} sur $R_{S'}$, et le foncteur
\struck{chgt de base} $R_{S'} \longrightarrow R_{S''}$ est \emph{continu}.

\marginal{on veut \ill{} que les foncteurs chgt. de base soient \emph{exacts} !}

\marginal{on veut que dans $\mathrm{Fl}_S$ \ill{} quotient \ill{}
\uncertain{déterminable}}

\marginal{dans $R_Y$, il \ill{} de \ill{} supposer $Y = e$ --- n'utilisant que
R1 et R3, R4}

\marginal{\ill{} que la famille des objets de $R_{S'}$ soit
\uncertain{loc. constante} finie}

\page{191}\note{sa page 2}

définit un morphisme de sites au sens inverse\,].

\emph{Axiomes d'une topologie} :
\begin{enumerate}
\item[a)] Si une famille est majorée par une famille couvrante, elle est
couvrante : trivial
\item[b)] Une somme est couvrante : trivial
\item[c)] stabilité par chgt. de base : on est ramené à des $X_i \to X$ qui
sont des monomorphismes, grâce à b). mais dans (\ill{}) $r$ commute aux
produits fibrés $X_i \times_X X'$, et on gagne.
\item[d)] Nature locale : supposons \ill{} qu'on ait des $X'_{\alpha} \to X$
couvrant $X$, tels que \struck{\ill{}} $\forall \alpha$, les
$X_i \times_X X'_{\alpha} = X'_{i\alpha}$ couvrent $X_{\alpha}$. Alors les
$X_i$ couvrent $X$. Ici encore, on est ramené au cas où les $X_i \to X$ sont
des monomorphismes, et on gagne.
\end{enumerate}

\begin{tikzcd}[column sep=small, row sep=small]
X_i & X'_i \arrow[l] \\
X \arrow[u] & X' \arrow[l] \arrow[u]
\end{tikzcd}

\begin{tikzcd}[column sep=small, row sep=small]
X_i \arrow[d] & X_{i\alpha} \arrow[d] \\
X & X_{\alpha} \arrow[l]
\end{tikzcd}

\ill{} SGA 4 : (Question : y a-t-il un axiome \ill{} l'image inverse d'une
topologie ? Mais attention, on \ill{} que ce n'est pas exact : g.\,!). \ill{}
on a utilisé \ill{} que R1. \emph{R3} \struck{\ill{}} \add{signifie} que les
topologies \ill{} $R_{S'}$ sont loc. \ill{} finies. \ill{} que ce canonique,
\emph{R4} que \ill{} L'axiome sur le changement de base est trivial.

\textbf{Corollaire.} Prenant les topos \struck{associés aux} $R_{S'}$, associé
aux $R_{S'}$, on trouve un \struck{\ill{}} \emph{champ en topos}
$\widetilde{R}_{S'}$ « sur » la sous-catégorie des objets loc. constants finis
de $R_{S'}$. Alors

\note{suit un bloc rayé de plusieurs diagonales, qui pose une question à
laquelle il ne répond pas : \struck{$X \to Y$ un morphisme dans $R_{S'}$,
$Y = \coprod_{n \in \mathbb{N}} Y_n$, avec $X \times_Y Y_n \to Y_n$ \ill{}
(somme effective) ; la famille $Y_n \to Y$ est-elle couvrante ? les
$r(Y_n) \to r(Y)$ comment \ill{} ??}}

\section*{Objets subordonnés, sous-champs $R^P$, et le cas galoisien (pages 192 à 194)}

\page{192}\note{sa page 3}

\marginal{Utilise R1 et R2}

\textbf{Prop. 2} Soit $X \in \mathrm{Ob}\,R_{S'}$. Alors
$X = \emptyset \Longleftrightarrow r(X) = \emptyset$, \; $X$ connexe
$\Longleftrightarrow r(X)$ connexe.

Soit $C$ une $\mathcal{U}$-\struck{catégorie} \emph{multigaloisienne}, $P$ un
objet. Un objet $X$ de $C$ est dit \emph{subordonné} à $P$ si
$X_P \to P$ est \struck{\ill{}} \ill{} $P$, i.e. $X_P$ est de \ill{} un, $X$
est \emph{connexe}.

\note{la marge porte en regard, d'une part un produit fibré de $X_P$ et d'un
second objet au-dessus de $X$, d'autre part le carré dont les sommets sont $X$,
ce second objet, $e$ et $P$ ; l'indice du second objet porte un trait que je ne
sais pas identifier, et il n'est pas conjecturé ici}

Pour $P$ fixé, les objets $X$ de $C$ subordonnés à $P$ forment une
\emph{sous-catégorie pleine} stable par $\varprojlim$ finies et par sommandes
directs

[car les objets \add{sur $P$} (qui sont « constants par morceaux ») forment une
sous-catégorie pleine de $C_{/P}$ stable par $\varprojlim$ finies et \ill{}
directes\,\ldots\,], donc c'est une \emph{catégorie multigaloisienne pleine} de
$C$, soit $C^{(P)}$. Si $P$ \struck{\ill{} de \ill{} $G$, cette catégorie
\ill{}} \ill{} maintenant $C = R_{S'}$, on peut regarder la

la sous-catégorie \emph{fibrée pleine} de $R$ formée des
$X' \in \mathrm{Ob}\,R_{S'}$, qui sont \emph{localement subordonnés} à $P$. On
voit encore que c'est \ill{} par $\varprojlim$ finies et par sommandes
directes, donc c'est une \struck{\ill{}} sous-catégorie
\emph{multigaloisienne pleine $R^P$ de $R$}.

\marginal{\ill{} de catégories multigaloisiennes \ill{} dans $P$}

Supposons maintenant $P$ \emph{galoisien} de groupe $G$, alors les objets de
$R_S$ \emph{subordonnés} (pas seulement loc.\,!) à $P$ correspondent aux objets
sur $P$ \emph{constants par morceaux}, avec \ill{} de $G$ dessus. Or, grâce à
R1, R\struck{2}\add{4}, un objet $X$ sur $P$ constant par morceaux
\add{$r(X)$} définit un objet \ill{} $r(P)$ constant par morceaux, et

\marginal{relèvement}

\page{193}

\note{page sans numéro de sa main, insérée entre ses pages 3 et 4 ; elle est
rayée de plusieurs longues diagonales}

\textbf{Proposition.} Supposons R1, R3, R4. Alors R2 équivaut à la condition
que $r$ est \emph{fidèle}.

Supposons $r$ fidèle, prouvons R2. Soient $X'$, $X''$ deux sous-objets de
$X \in \mathrm{Ob}\,R_{S'}$, tels que $r(X') \supset r(X'')$ comme sous-objets
de $r(X)$, \struck{\ill{}} or
\[
  r(X' \cap X'') = r(X') \cap r(X''),
\]
(\struck{def} lemme) \struck{\ill{}} les inclusions
$X' \cap X'' \subset X'$ et $X' \cap X'' \subset X''$ \struck{\ill{}} \ill{}
isom. dans $r$. On est ramené au cas où on a \add{déjà} une inclusion
$X' \subset X''$, i.e. $X'' = X' \sqcup X_1$. Alors
$r(X'') = r(X') \sqcup r(X_1)$ et l'hypothèse $r(X') = r(X'')$ équivaut à
$r(X_1) = \emptyset$.

\page{194}\note{sa page 4}

\struck{\ill{}} utilisant de plus R2, on voit que le foncteur
$X \mapsto r(X)$ des \struck{\ill{}} \emph{objets constants par morceaux} dans
les \ill{} sur $r(P)$ constants par morceaux est \emph{pl. fidèle}. Utilisant
de plus R2 et R3, on trouve qu'il est \uncertain{essentiellement surjectif}.
C'est donc une \emph{équivalence} de catégories, \struck{i.e.} Notons que comme
$G$ opère sur $R$ il opère sur $r(P)$, \struck{donc} et on trouve que
l'équivalence est compatible avec les opérations de $G$. On trouve ainsi :

\marginal{lemme déjà \ill{} utilisé \ill{}}

\textbf{Prop. 3} (Soit $P$ un objet \emph{galoisien de groupe $G$ dans
$R_{S'}$}.) Le foncteur $X \mapsto r(X \times P)/r(P)$ \ill{} une
\emph{équivalence} de la catégorie des objets de $R_{S'}$ subordonnés à $P$,
avec la catégorie des \struck{\ill{} étales} (constants par morceaux sur
$r(P)$, sur lesquels $G$ opère de façon compatible avec ses opérations sur
$r(P)$).

Utilisant le théorème de la descente, il vient

\textbf{Cor. 1} Soit $P$ objet \emph{galoisien de groupe $G$ dans $R_S$}. Le
foncteur $X \mapsto r(X \times P)/r(P)$ induit une \emph{équivalence} de la
catégorie des objets de $R_S$ loc. \emph{subordonnés} à $P$, avec la catégorie
des \struck{revêtements étales} \struck{\ill{}} de $Z = r(P)$, \emph{à action
de $G$ compatible} avec ses actions sur $r(P)$, et satisfaisant la condition
suivante (automatiquement vérifiée si $r(P) = Z$ est fini sur $S$) : les
$S' \to S$ étales \emph{tels que} $Z'_{S'}$ \struck{soient étales sur} soient
constants par morceaux \ill{} $Z_{S'}$ recouvrent $S$.

\note{Localisant sur $S$, on trouve $=$ une description \struck{du} champ $R^P$
en termes des schémas $Z$ à groupe d'opérateurs $G$ sur $S$ --- la remarque est
de lui, écrite à la suite après un trait oblique}

\section*{Le champ associé à un schéma à opérateurs, et pourquoi il faut recoller (pages 195 et 196)}

\page{195}\note{sa page 5}

Notons d'ailleurs

\textbf{Prop. 4} Pour tt schéma \struck{\ill{}} $Z$ sur $S$, à groupe fini $G$
d'opérateurs, \struck{\ill{}} agissant de \emph{façon admissible}, le champ
\struck{défini} \add{décrit} \struck{par le} \add{comme} dessus est un champ en
\emph{$\mathcal{U}$-catégories multigaloisiennes}, et le foncteur « passage au
quotient par $G$ » (qui est défini, car l'action de $G$ est \emph{admissible})
\struck{est} \ill{} une \uncertain{donnée} de ramification sur $S$.

\note{le mot que je lis \emph{donnée} pourrait être \emph{domaine} : c'est le
terme que la page 190 définit et que la chemise de la page 188 porte en titre ;
la lettre ne tranche pas et l'article, féminin, va contre}

Que $R_{S'}$ soit multigaloisienne : tout d'abord, il \ill{} admis de la
catégorie des \struck{objets} \emph{rev. à opérateurs $G$} \struck{des}
$Z_{S'}$, où on en \ill{} prend \ill{} une \emph{sous-catégorie stable par
$\varprojlim$ finies et par sommandes directs} ; OK. \struck{Foncteur} \ill{}
et \emph{\ill{}} : OK.

Le foncteur $r$, étant un \emph{passage au quotient}, commute aux
$\varinjlim$ finies (que \ill{} calcule au sens de $(\mathrm{Sch})$\,!), d'où
\textbf{R1}. De $=$ \textbf{R2} est immédiat. \textbf{R3} aussi. \textbf{R4}
aussi. OK.

\struck{Ainsi} Supposons pour simplifier \ill{} $S$ \ill{} affine, et que $R$
satisfait une \emph{petite condition supplémentaire} (\ill{} vérifiée en
pratique). \struck{\ill{}} Pour tt objet de $R_S$, son \ill{} sur \ill{} de
$R_S$ \struck{\ill{}} base, i.e. \struck{$X$ est} \emph{subordonné} à un
\struck{objet} \emph{objet à opérateurs galoisien $(P, G)$ dans $R_S$}.

Alors $R_S$ est \emph{limite inductive filtrante} de \ill{} sous-catégories
\struck{\ill{}} \ill{} $R^P_S$. On \ill{} donne une \ill{} bonne description de
$R_S$, \ill{}, en termes d'un \emph{système projectif} \struck{projectif}
d'objets à opérateurs $(Z_i, G_i)$

\page{196}\note{sa page 6}

correspondants aux $(P_i, G_i)$.

Dans le cas où l'objet \emph{final} de $R_S$ est \emph{connexe}, on peut
prendre les $P_i$ \emph{connexes}, alors les morphismes de transition
$G_j \to G_i$ sont \emph{surjectifs}, et on trouve
$Z_i \simeq Z_j / H_{ij}$\,\ldots\ D'ailleurs, les $Z_i/G_i$ seront
connexes\,\ldots\ D'autre part, relativement à un tel système projectif, on
peut vérifier un énoncé comme précéd\add{emmen}t, aux

\textbf{Prop. 5} \ldots

Mais \struck{\ill{}} \emph{partant de} $R$, et \emph{construire} un
\struck{\ill{}} objet $(R_i, G_i)$, \ill{} \ill{} un
\struck{nouveau champ} $(R', r')$, et $R \hookrightarrow R^{\bullet}$, il n'est
pas vrai en général, \add{que} $R' \hookrightarrow R$ soit une équivalence :
\emph{Le champ $R$ n'est pas nécessairement « engendré » par ses sections
globales.} Donc pour définir un $R$ en \emph{général}, il faut \struck{y}
localiser et \emph{recoller}\,\ldots\ On va donc décrire un procédé de
\emph{recollement} qui nous permettra de construire les domaines de
ramification dont on aurait besoin.

\emph{Notions auxiliaires} : Si $(R, r)$, $(R', r')$ sont deux domaines de ram.
sur $S$, un \emph{hom.} de l'un dans l'autre est un \emph{couple
$(\varphi, \lambda)$} \add{\struck{avec fibres}} où
$\varphi : R \longrightarrow R'$ est un \emph{foncteur continu}, exact, et
$\lambda$ un \emph{isomorphisme} de foncteurs continus
$r'\varphi \simeq \struck{\ill{}}\, r$. \struck{\ill{}} De cette façon,

\section*{La 2-catégorie $\mathrm{Domram}(S)$, et le recollement (pages 197 et 198)}

\page{197}\note{sa page 7 ; le coin porte en outre la notation
$\mathrm{Domram}(S)$}

les domaines de ramification \add{du \ill{}} sur $S$ forment une
\emph{2-catégorie}. Pour $S'$ variable, les 2-catégories ainsi obtenues
$\mathrm{Domram}(S')$ forment une \emph{2-catégorie fibrée}\,\ldots

Si $(R, r)$ est défini par $Z$, $G$, de sorte qu'on \ill{} un objet
\emph{galoisien canonique} $P$ dans $R_S$, de groupe $G$, un
\struck{\ill{}} hom. $(R, r) \to (R', r')$ définit un objet $\varphi(P)$ de
$R'$, qui est \emph{galoisien de groupe $G$}, et un \emph{isomorphisme
compatible avec $G$}
\[
  r(\varphi(P)) \overset{\alpha}{\simeq} \struck{\ill{}}\; Z
\]

\textbf{Prop. 6} \struck{On} donne ainsi une \emph{équivalence} entre
$\underline{\mathrm{Hom}}((R, r), (R', r'))$ et la catégorie des \emph{couples}
$(P', \alpha)$, i.e. $P'$ est un objet \emph{galoisien} de $R'_S$, de groupe
$G$, et $\alpha$ un \emph{$G$-isomorphisme}
\[
  r(P') \simeq Z .
\]

[ Si $(P', \alpha)$ \struck{sont} deux termes, on a une \emph{équivalence}
$R \cong R'^{P'}$. On trouve en passant \struck{\ill{}} :

\textbf{Cor.} \struck{Alors} \add{Tout} hom. $(\varphi, \lambda)$ de domaines
de ramification, \struck{\ill{}} $\varphi$ est néc. \emph{pleinement fidèle}. ]

En particulier

\textbf{Cor.} \struck{En} Si $(R', r')$ est défini par $Z'$, $G'$, alors
$\underline{\mathrm{Hom}}((R, r), (R', r'))$ est \emph{équivalent} à la
catégorie des \struck{revêtements} \emph{principaux} $P'$ de $Z'$, de groupe
$G$, sur lequel $G'$ opère de façon compatible avec ses opérations sur $Z'$,
\struck{\ill{}} \add{muni} d'un isom. de $G$-schémas sur $S$
\[
  P'/G' \simeq Z .
\]

\page{198}\note{sa page 8}

On explicite \struck{facilement} la \emph{composition} des hom. de domaines de
ramification associés : $Z, G$ ; $Z', G'$ ; $Z'', G''$ \ldots\ Ceci dit :

Soient des $S'_i$ couvrant $S$, sur chaque $S'_i$ \emph{étant} donné un
$(Z_i, G_i)$ \emph{schéma à groupe d'opérateurs $G_i$ admissible}, pour tt
couple $(i,j)$, on suppose donné \struck{\ill{}} sur $S_i \times_S S_j$ une
« \emph{subordination} » de $(Z_{ij}, G_i)$ à $(Z_{ji}, G_j)$
$\bigl(Z_{ij} = Z_i \times_{S_i} (S_i \times_S S_j)\bigr)$ et
\struck{\ill{}} sur les triples \ill{} des isomorphismes de transitivité entre
les subordinations. On suppose enfin une \emph{condition de compatibilité sur
les triples}. Alors on peut \emph{recoller} les $(R_i, r_i)$ définis par les
$(R_i, G_i)$ \ldots

Donnons-en en particulier \ill{} \ill{} de \ill{} le cas de nous :

Soit une catégorie \emph{fibrée} \add{(de « modèles »)} $C$ sur
$\mathrm{Et}(S)$ (pas nécessairement un champ) et un \emph{foncteur} continu
$s : C \longrightarrow \mathrm{Fl}_S$. On suppose $C$ fibrée en
\emph{groupoïdes}, que deux objets d'un $C_{S'}$ sont \emph{localement
isomorphes}, que $C_{S'}$ est \struck{localement} non vide --- d'où, en \ill{}
\ill{} en \emph{champ associé}, une \emph{gerbe}. On suppose que le
\emph{lien} de cette gerbe est défini \struck{\ill{}} \emph{loc. constant}
\ill{}, i.e. pour tt $x \in C_{S'}$, le \emph{faisceau associé}

\page{199}\note{sa page 9}

aux préfaisceaux $S'' \longrightarrow \mathrm{Aut}\, x_{S''}$ est
\emph{représentable} par un \struck{\ill{} la catégorie des} \emph{rev. ét.} de
$S'$. On a une « \emph{représentation} » de cette \emph{gerbe} dans le
\emph{préchamp} $\mathrm{Fl}_S$. À un tel lien, on associe un \emph{domaine de
ramification} sur $S$ ainsi :

\begin{tikzcd}[column sep=small, row sep=small]
& \mathrm{Rev}\,\mathrm{Fl}_S \arrow[d] \\
\mathcal{G} \arrow[r] \arrow[d] & \mathrm{Fl}_S \\
\mathrm{Et}(S) &
\end{tikzcd}

les objets de $R_{S'}$ sont les hom. continus
$\mathcal{G} \longrightarrow \mathrm{Rev}\,\mathrm{Fl}_S$ qui relèvent
$\mathcal{G} \to \mathrm{Fl}_S$, et item quand $r$ est remplacé par $r'$.

Le foncteur $r_S : R_S \longrightarrow \mathrm{Et}_{/S}$ est un
foncteur \uncertain{lien}. On vérifie que tt marche.

\note{le mot souligné en fin de phrase porte un trait qui le fait ressembler
aussi bien à un $\varinjlim$ qu'à \emph{lien} ; le sens de la phrase n'est pas
fixé par la page}

\textbf{Ex. 1} Donnons \struck{\ill{}} une famille
$\underline{a} = (a_i)_{i \in I}$ de sections de $\mathcal{O}_S$, premières
avec \ill{}, et une famille $\underline{n} = (n_i)_{i \in I}$ d'entiers
$\geq 1$, \struck{\ill{}} Soit \ill{} un ens. fini, \ill{} par les groupes
$C_{S'} = \struck{\ill{}}$ groupoïde \ill{} (une gerbe dont \ill{}) et
\[
  s(C_{S'}) = Z^{\underline{a}}_{\underline{n}} , \qquad
  \mu_{\underline{n}}(S') = \prod_{i \in I} \mu_{n_i}(S) ,
\]
\ill{} opération qui en dérive de ce groupe. Les axiomes sont vérifiés. [ En
fait, on \ill{} comporte en cas d'une \emph{gerbe triviale} de \ill{}
$Z^{\underline{a}}_{\underline{n}}$ sur lequel $\mu_{\underline{n}}$ opère ].

\textbf{Ex. 1 bis} Donnons \struck{\ill{}} \ill{} $\underline{a}$,
$\underline{n}$ \ill{} : le \emph{cas commutatif}, mais on \ill{} : la limite
sur $\underline{n}$ de \ill{}, $=$ le \emph{domaine de ramification
kummerienne} défini par $\underline{a}$.

\textbf{Ex. 2} Donnons \add{une} famille $(D_i)_{i \in I}$ de \emph{diviseurs},
et une famille $(n_i)_{i \in I}$ d'\ill{} \ill{} faisant \ill{} \ill{}
construction que dans l'ex. 1. On prend

\page{200}\note{sa page 10}

$C_{S'}$ défini ainsi :
\[
  \mathrm{Ob}(C_{S'}) = \bigl\{\, \underline{a} = (a_i)_{i \in I}
  \;\bigm|\; \forall i,\; a_i \text{ une équation de } D_i \,\bigr\}
\]
\[
  \mathrm{Hom}(\underline{a}, \underline{b}) = \bigl\{\, \xi = (q_i)_{i \in I}
  \;\bigm|\; \forall i,\; b_i = a_i \xi_i^{n_i} \,\bigr\}
\]
\struck{Isomorphismes et} \emph{Composition} des lois \emph{évidentes}. Images
inverses \ill{}. On trouve ainsi \ill{} champ associé une \emph{gerbe
abélienne} de groupe $\mu_{\underline{n}}$. (L'existence d'une \emph{section
dépend} de la \emph{nullité} d'un élément de $H^2(X, \mu_{\underline{n}})$, qui
n'est autre, comme on devine, que \struck{$C$} la \emph{classe de cohomologie
des $D_i$ dans les $H^1(X, \mu_{n_i})$}\,!\,) : \emph{domaine de ramification
kummerienne} défini par $(D_i)_{i \in I} = D$ et $\underline{n}$.

\note{les deux exposants de $H$ sont tracés de la même main qui ne distingue
pas le $1$ du chiffre romain, et le second pourrait être un $I$ ; ils sont
laissés tels qu'ils sont écrits}

\textbf{Ex. 2 bis} On donne $D = (D_i)_{i \in I}$, on pose : la
\emph{limite sur $\underline{n}$} : \emph{domaine de ramification kummerienne
défini par $D$}.

Ces exemples boitent un tant \add{soit} peu, on voudrait au fait que \ill{} ne
sont \add{pas} premiers avec les \uncertain{résiduelles} des $D_i$ seulement,
\ill{} pour \ill{}. \ill{} dans le cas de l'exemple 1, on \ill{} de tt $S$.
Dans le cas de l'exemple 1, on prendra les \emph{rev. étales} de
$Z^{\underline{a}}_{\underline{n}}$, sur lesquels $\mu_{\underline{n}}$ opère
de \emph{façon compatible avec ses opérations}, sur
$Z^{\underline{a}}_{\underline{n}}$. Dans le cas de l'exemple 2

\note{la page s'arrête là, en cours de phrase ; la suite passe dans le lot
suivant}

\end{document}
